\documentclass[UTF8,a4paper,12pt]{ctexart}

\usepackage[left=2.5cm,right=2.5cm,top=2.5cm,bottom=2.5cm]{geometry}
\usepackage{setspace}
\usepackage{amsmath,amssymb,amsthm}
\usepackage{booktabs}
\usepackage{longtable}
\usepackage{array}
\usepackage{graphicx}
\usepackage{float}
\usepackage{enumitem}
\usepackage{hyperref}
\usepackage{caption}

\hypersetup{colorlinks=true,linkcolor=black,citecolor=black,urlcolor=blue}
\setstretch{1.25}
\graphicspath{{../results/figures/}{figures/}}

\title{第十八届\u201c华中杯\u201d大学生数学建模挑战赛A题\\论文初稿（LaTeX 版）}
\author{参赛队号：\underline{\hspace{4cm}}}
\date{\today}

\begin{document}

\maketitle

\begin{abstract}
针对城市绿色物流配送调度问题，本文建立了\u201c静态基础模型 + 政策约束扩展 + 动态重调度机制\u201d的分层优化框架。问题一构建带软时间窗、双容量约束与时变速度的混合车队车辆路径模型，并以固定成本、运营成本和碳排放成本之和最小为目标；问题二在基础模型上加入绿色配送区限行约束，分析政策实施前后的成本与排放变化；问题三针对订单取消、新增、地址变更和时间窗调整等突发事件，采用事件触发的滚动时域重规划策略实现快速可行修复。求解方法采用\u201c构造启发式 + 局部搜索\u201d组合算法，并通过多维指标评价方案有效性。本文给出的论文框架与公式体系可直接用于后续数据实验、图表填充与定稿排版。

\textbf{关键词：}绿色物流；车辆路径问题；软时间窗；混合车队；动态重调度
\end{abstract}

\tableofcontents
\newpage

\section{问题重述与分析}

\subsection{问题背景}
城市配送系统在双碳目标与城市交通治理政策的共同驱动下，面临\u201c成本控制、时效保障、低碳约束\u201d三重压力。题目聚焦于多车型联合调度场景，要求在满足客户服务质量的同时，优化配送路径与车辆配置。

\subsection{问题分解}
结合题意，可将任务拆解为三类子问题：
\begin{enumerate}[label=\arabic*)]
    \item 静态场景下的混合车队路径优化；
    \item 绿色配送区限行政策加入后的约束扩展与对比分析；
    \item 订单动态变化场景下的实时重调度策略设计。
\end{enumerate}

\subsection{总体建模思路}
采用分层建模策略：先构建问题一的基础模型，再通过增量约束扩展处理问题二，最后在滚动时域框架下处理问题三的动态事件。该策略可降低模型耦合复杂度，提升工程实现稳定性与可解释性。

\section{模型假设与符号说明}

\subsection{模型假设}
为保证模型可解并与题目信息一致，作如下假设：
\begin{enumerate}[label=\arabic*)]
    \item 所有车辆均从同一配送中心出发并最终返回配送中心；
    \item 路段速度由时段状态决定，可用分段期望速度近似；
    \item 客户服务时长固定（可取 $20$ 分钟）；
    \item 充电、加油等待时间在基础模型中忽略；
    \item 时间窗采用软约束，违反部分通过惩罚项计入目标函数。
\end{enumerate}

\subsection{符号定义}
\begin{table}[H]
\centering
\caption{核心符号说明}
\begin{tabular}{p{0.2\textwidth}p{0.65\textwidth}}
\toprule
符号 & 含义 \\
\midrule
$N=\{0,1,\dots,n\}$ & 节点集合，$0$ 为配送中心 \\
$K$ & 车辆集合 \\
$T$ & 车型集合 \\
$d_{ij}$ & 节点 $i$ 到节点 $j$ 的距离 \\
$q_i^w,q_i^v$ & 客户 $i$ 的重量、体积需求 \\
$[e_i,l_i]$ & 客户 $i$ 的服务时间窗 \\
$Q_t^w,Q_t^v$ & 车型 $t$ 的载重与容积上限 \\
$x_{ijk}^t$ & 是否由车型 $t$ 的车辆 $k$ 执行弧 $(i,j)$ \\
$y_{ik}^t$ & 客户 $i$ 是否由车型 $t$ 的车辆 $k$ 服务 \\
$t_i$ & 到达客户 $i$ 的时刻 \\
$z_k^t$ & 车型 $t$ 的车辆 $k$ 是否启用 \\
\bottomrule
\end{tabular}
\end{table}

\section{问题一：静态基础模型}

\subsection{目标函数}
以总成本最小为优化目标，综合考虑车辆固定启用成本、运输能耗成本、碳排放成本与时间窗惩罚：
\begin{equation}
\min Z = \sum_{t\in T}\sum_{k\in K} C_{\text{fix}} z_k^t
+ \sum_{t\in T}\sum_{k\in K}\sum_{i\in N}\sum_{j\in N}
\left(C_{\text{op}}^{ijk,t} + C_{\text{em}}^{ijk,t}\right)
+ \sum_{i\in N} P_i .
\end{equation}

其中软时间窗惩罚可写为
\begin{equation}
P_i = \alpha\max(e_i-t_i,0) + \beta\max(t_i-l_i,0).
\end{equation}

\subsection{约束条件}
\textbf{(1) 双容量约束}
\begin{align}
\sum_{i\in N} q_i^w y_{ik}^t &\le Q_t^w,\quad \forall k\in K,\forall t\in T,\\
\sum_{i\in N} q_i^v y_{ik}^t &\le Q_t^v,\quad \forall k\in K,\forall t\in T.
\end{align}

\textbf{(2) 路径连续性约束}
\begin{align}
\sum_{j\in N, j\neq i} x_{ijk}^t &= y_{ik}^t, \quad \forall i,k,t,\\
\sum_{i\in N, i\neq j} x_{ijk}^t &= y_{jk}^t, \quad \forall j,k,t.
\end{align}

\textbf{(3) 时间递推约束}
\begin{equation}
t_j \ge t_i + s_i + \frac{d_{ij}}{v_{ij}} - M(1-x_{ijk}^t),\quad \forall i,j,k,t.
\end{equation}

\textbf{(4) 车辆可用数量约束}
\begin{equation}
\sum_{k\in K} z_k^t \le N_t,\quad \forall t\in T.
\end{equation}

\textbf{(5) 变量取值约束}
\begin{equation}
x_{ijk}^t,y_{ik}^t,z_k^t\in\{0,1\},\quad t_i\ge 0.
\end{equation}

\section{问题二：绿色配送区政策约束扩展}

\subsection{约束扩展思路}
问题二不改变问题一的主体结构，仅在基础模型上增加\u201c区域-时段-车型\u201d耦合约束。设 $g_i$ 表示客户 $i$ 是否位于绿色配送区，$\tau_i$ 表示服务时刻。

当客户位于绿色区且服务时刻处于限行时段（如 $[8{:}00,16{:}00]$）时，燃油车型不允许服务：
\begin{equation}
g_i=1,\ \tau_i\in[8,16] \Rightarrow y_{ik}^t=0,\quad \forall t\in T_{\text{fuel}},\forall k\in K.
\end{equation}

在实现中可将其改写为线性不等式，便于求解器处理。

\subsection{对比评价指标}
为体现政策效应，建议比较如下指标：
\begin{enumerate}[label=\arabic*)]
    \item 总成本变化率；
    \item 总碳排放变化率；
    \item 燃油车与新能源车的任务占比变化；
    \item 平均客户迟到时长变化。
\end{enumerate}

\subsection{结果展示建议}
图表建议至少包括：
\begin{enumerate}[label=\arabic*)]
    \item 政策前后车辆类型使用对比柱状图；
    \item 政策前后路径示意图；
    \item 成本与排放对比表。
\end{enumerate}

\section{问题三：动态事件驱动重调度模型}

\subsection{动态事件类型}
考虑四类常见事件：订单取消、新增订单、客户地址变化、时间窗更新。问题三目标不是每次求全局最优，而是在有限计算时间内快速生成高质量可行解。

\subsection{滚动时域机制}
设重规划时刻为 $t^r$，滚动窗口长度为 $H$。在每个事件触发时刻，仅对未完成任务子集进行局部重优化：
\begin{equation}
\min \ Z\bigl(\mathcal{R}(t^r,H)\bigr)
\end{equation}
其中 $\mathcal{R}(t^r,H)$ 表示当前滚动窗口内需要重新决策的任务集合。

\subsection{快速修复策略}
采用\u201c插入-删除-交换\u201d三类算子进行局部修复：
\begin{enumerate}[label=\arabic*)]
    \item 新增订单：最小增量成本插入；
    \item 订单取消：从既有路径删除并局部平衡；
    \item 地址/时间窗变化：邻域内重排并二次可行性检查。
\end{enumerate}

\subsection{算法终止准则}
可使用以下任一条件作为终止准则：
\begin{enumerate}[label=\arabic*)]
    \item 达到最大迭代次数；
    \item 连续若干次迭代无改进；
    \item 达到实时响应时限（如 30 秒）。
\end{enumerate}

\section{求解流程与结果分析（待填充）}

\subsection{求解算法框架}
建议采用\u201c构造启发式 + 局部搜索\u201d两阶段流程：
\begin{enumerate}[label=\arabic*)]
    \item 初始解构造：聚类/扫描法 + 最近邻插入；
    \item 邻域改进：2-opt、跨路径交换、重定位；
    \item 约束修复：时间窗与容量违规惩罚回退；
    \item 输出结果：路径、成本、排放、服务质量指标。
\end{enumerate}

\subsection{结果表格模板}
\begin{table}[H]
\centering
\caption{三类问题结果汇总（示例模板）}
\begin{tabular}{cccccc}
\toprule
场景 & 总成本(元) & 碳排放(kg) & 车辆数(辆) & 平均迟到(h) & 求解时间(s) \\
\midrule
问题一 & 待填 & 待填 & 待填 & 待填 & 待填 \\
问题二 & 待填 & 待填 & 待填 & 待填 & 待填 \\
问题三 & 待填 & 待填 & 待填 & 待填 & 待填 \\
\bottomrule
\end{tabular}
\end{table}

\subsection{图件插入模板}
\begin{figure}[H]
    \centering
    \includegraphics[width=0.8\textwidth]{route_example.png}
    \caption{路径规划结果示意图（占位图，待替换）}
\end{figure}

\subsection{敏感性分析建议}
建议至少分析以下敏感性：
\begin{enumerate}[label=\arabic*)]
    \item 碳排放成本系数变化对车型选择的影响；
    \item 迟到惩罚系数变化对时间窗满足率的影响；
    \item 绿色区半径或限行时段变化对总成本的影响。
\end{enumerate}

\section{模型评价与结论}

\subsection{模型优点}
\begin{enumerate}[label=\arabic*)]
    \item 分层建模结构清晰，便于逐问扩展；
    \item 双容量、软时间窗、政策约束与动态事件均可统一纳入；
    \item 算法工程实现成本低，适合比赛时限内快速迭代。
\end{enumerate}

\subsection{模型局限}
\begin{enumerate}[label=\arabic*)]
    \item 时变速度采用分段近似，未显式考虑随机波动传播；
    \item 动态重调度以局部最优为主，缺乏全局回溯机制；
    \item 未纳入充电排队、司机工时等更细粒度运营约束。
\end{enumerate}

\subsection{结论与后续工作}
本文完成了面向 A 题的论文初稿结构与核心模型搭建。后续可基于程序输出结果补充图表、完成参数标定与消融实验，并进一步优化语言表达与排版细节以形成终稿。


\bibliographystyle{unsrt}
\bibliography{refs}

\appendix
\section{附录：符号补充表}
\begin{longtable}{p{0.2\textwidth}p{0.55\textwidth}p{0.15\textwidth}}
\toprule
符号 & 含义 & 单位 \\
\midrule
\endfirsthead
\toprule
符号 & 含义 & 单位 \\
\midrule
\endhead
$N$ & 客户点集合，$0$ 表示配送中心 & - \\
$K$ & 车辆集合 & - \\
$T$ & 车型集合 & - \\
$d_{ij}$ & 节点 $i$ 到节点 $j$ 的道路距离 & km \\
$q_i^w$ & 客户 $i$ 的重量需求 & kg \\
$q_i^v$ & 客户 $i$ 的体积需求 & m$^3$ \\
$[e_i,l_i]$ & 客户 $i$ 的时间窗 & h \\
$s_i$ & 客户 $i$ 的服务时间 & h \\
$Q_t^w,Q_t^v$ & 车型 $t$ 的最大载重与最大容积 & kg, m$^3$ \\
$x_{ijk}^t$ & 车型 $t$ 的车辆 $k$ 是否由 $i$ 行至 $j$ & 0-1 \\
$y_{ik}^t$ & 客户 $i$ 是否由车型 $t$ 的车辆 $k$ 服务 & 0-1 \\
$t_i$ & 车辆到达客户 $i$ 的时刻 & h \\
$z_k^t$ & 车型 $t$ 的车辆 $k$ 是否被启用 & 0-1 \\
\bottomrule
\end{longtable}

\end{document}
